% generated from src/5 Conclusion.md by tools/doc2tex.py; do not edit
\chapter{CONCLUSION AND FUTURE WORK}

\section{Answers to the research questions}

\textbf{RQ1, feasibility at the edge: yes.} On a Raspberry Pi 5 with a Hailo-8, the deployment setting (n model with tiles) ran 15 minutes in real time:

\begin{itemize}
\tightlist
\item
  30 FPS with 45 ms latency
\item
  4.3 W measured on the board, plus about 2.5 W for the accelerator
\item
  peak SoC temperature 60 °C, no throttling
\end{itemize}

The pipeline is limited by the accelerator, not the CPU. The m model with tiles reaches only 7.9 FPS, enough for analysis but not for a live camera without dropping frames.

\textbf{RQ2, coverage: substantially improvable without retraining.} Tiling the back of the room raised the share of SCB students found from 49\% to 76\%, and in the far third from 17\% to 74\%. It doubled the students found in a dense study hall. Removing same-skeleton duplicates and anchoring seats at the head made the inferred class size match the real one (52 for about 50 students, 31 for about 30).

In a small room filmed at close range (clip\_5), seats still multiplied, to 13-17 for 10 students, so the class-size estimate holds for the camera distances it was developed on, not yet for very close cameras.

\textbf{RQ3, the per-student decision: calibration beats benchmark accuracy.} Classifiers trained on SCB rank SCB actions well (within-image AUC 0.84) but transfer badly:

\begin{itemize}
\tightlist
\item
  \textbf{Lecture clip:} they call 48-89\% of an attentive lecture distracted.
\item
  \textbf{Human labels:} they wrongly flag 114-199 of 262 focused moments.
\item
  \textbf{Cause:} SCB has no ``listening'' label, so they learn that a raised head means distracted.
\end{itemize}

The calibrated rules have very few false alarms (specificity 90\% on human labels; 0-5\% on the lecture). They first missed most real distraction (11\% on the staged session), because the ear-based head yaw collapses in profile. A second, body-relative head-turn measure raised that to 45-56\%, with no new false alarms there and a small cost on real classes. On a 10-student class labelled after all thresholds were fixed (clip\_5), the rules on the deployment setting kept 91\% specificity. Phone use remains invisible to pose.

\textbf{RQ4, the class indicator: yes, with honest width.} Seat memory, zone weighting and a Wilson interval around the shown value give a share that always lies inside its range. Without the real head count the range is wide (25-37 points for 20-30 students), which reflects how little a sample that small can say. Giving \texttt{-\/-class-size} narrows it.

\section{Future work}

\begin{enumerate}
\def\labelenumi{\arabic{enumi}.}
\tightlist
\item
  \textbf{A held-out test at classroom scale.} Run the 9-person scripted session (research/staged/script\_v2.csv). clip\_5 confirmed the rules' specificity on a class nothing was tuned on, but its recall still rests on one two-person session.
\item
  \textbf{Classroom-specific training labels.} The labelling and staged-session tools can grow a dataset that includes attentive, head-up students, the case SCB lacks. A learned decision trained on it, with per-seat baselines as features, could combine the classifiers' sensitivity with the rules' specificity.
\item
  \textbf{Phone use.} Detecting phones needs an object detector. On the Hailo-8 it could run on the same frames, at a cost in frame rate. Only a per-seat ``phone visible'' flag would leave the device.
\item
  \textbf{Longer field deployments.} Weeks of use would show drift in calibration, the effect of lighting and seating changes, and whether teachers find the indicator useful. That needs a study with consent, designed with the school.
\item
  \textbf{Live m with tiles.} A leaky queue or a lower camera rate would let the most accurate setting run live without growing latency.
\item
  \textbf{Keypoint-based geometry.} Anchor seats on the shoulders and measure torso rotation against the visible upper body instead of the box. That removes the two failures seen at close range in clip\_5: seats multiplying and a full-body box reading as a turned torso.
\item
  \textbf{Teacher-facing evaluation.} The real test of a class-level indicator is whether it helps teaching, which needs a user study rather than more accuracy metrics.
\end{enumerate}
